% !TEX root = TI5.tex

\section{Appendix for proofs that are very similar to prior work} \label{app:similar}
In this appendix, we gather together all the proofs of claimed results where our result generalizes a prior result, but the proof is very nearly the same as the original result, one must merely check carefully that the corresponding statements in our generalized setting remain true. In each proof, while we include much of the proof to check that the result still holds in our generalized setting, we also highlight any differences needed from the original proof in \textbf{bold}.

\begin{proof}[Proof of Theorem~\ref{thm:HillarRheaPIR}]
Following this notation, we can think of elements of $M = \bigoplus_{i=1}^n R/\m^{e_i}$ as column vectors with $n$ entries, where the $i$-th entry is an element of $R/\m^{e_i}$. Let $\iota_i \colon R/\m^{e_i} \to M$ be the natural inclusion map of the $i$-th summand, and let $\pi_j \colon M \to R/\m^{e_j}$ be the natural quotient map with kernel $\bigoplus_{i \neq j} R/\m^{e_i}$, such that $\pi_i \iota_i$ is the identity on $R/\m^{e_i}$ for each $i$. For $m \in M$, let $m_i = \pi_i(m)$, then $m$ can be written as $m = \iota_1(m_1) + \dotsb + \iota_n(m_n)$, or in the more usual vector notation as $(m_1, \dotsc, m_n)^T$ (the transpose just because we will find it convenient to think of these as column vectors). 

For an endomorphism $\varphi \colon M \to M$, let $\varphi_{ij} = \pi_i \circ \varphi \circ \iota_j$, the latter being an element of $\Hom(R/\m^{e_j}, R/\m^{e_i})$. If $\psi \in \End(M)$, let $\sigma = \varphi + \psi$ and $\mu = \psi \varphi$, and define $\psi_{ij}, \sigma_{ij}, \mu_{ij}$ analogously. Then it is routine to verify that $\sigma_{ij} = \varphi_{ij} + \psi_{ij}$ and $\mu_{ij} = \sum_{k=1}^n \psi_{ik} \circ \varphi_{kj}$. Thus this set of matrices forms a ring in the usual manner, and the map taking $\varphi$ to the matrix whose $(i,j)$ entry is $\varphi_{ij}$ for all $i,j \in \{1,\dotsc,n\}$ is a ring homomorphism from $\End(M)$ into this ring of matrices. It is routine to verify that this homomorphism is injective and surjective, and is thus an isomorphism. This completes the proof of the first part of the result.

Next, consider any $R$-module homomorphism $\varphi \colon R/\m^{e_j} \to R/\m^{e_i}$. Because these are cyclic $R$-modules, such a homomorphism is completely determined by $\varphi(1)$, for we have $\varphi(r) = r \cdot \varphi(1)$ for all $r \in R$. Thus $\varphi$ is completely specified by the single element $\varphi(1) \in R/\m^{e_i}$, and the only constraint on the latter is that it must be the case that multiplication by $\m^{e_j}$ annihilates it, since that was true of $1$ in $R/\m^{e_j}$. Now, if $e_i \leq e_j$, this is true of every element of $R/\m^{e_i}$, and there is no constraint on $\varphi(1)$. If $e_i > e_j$, then $\varphi(1)$ must be in $\m^{e_i - e_j}$, or equivalently, must be divisible by $t^{e_i - e_j}$, where $t$ is a generator of $\m$.

Combining our description of elements of $\Hom(R/\m^{e_j}, R/\m^{e_i})$ with our description of endomorphisms of $M$, we see that endomorphisms of $M$ can equivalently be thought of as $n \times n$ matrices $A$ where the elements in the $i$-th row are elements of $R/\m^{e_i}$, and such that $t^{e_i - e_j}$ divides $A_{ij}$ whenever $e_i \geq e_j$. Let $\overline{S}$ denote the preceding set of matrices, which we have shown is isomorphic to $\End(M)$.

Finally, it remains to characterize which matrices in $\overline{S}$ are invertible. Suppose $A,B \in \overline{S}$ satisfy $AB = I_n$, the $n \times n$ identity matrix. Taking the latter equation modulo $\m$, we find that the matrix $\overline{A}$ gotten by taking every entry of $A$ modulo $\m$ must be invertible. All the remains is to show the converse.

So suppose $A$ is an element of $\overline{S}$, such that $A \pmod{\m}$ is invertible as a matrix over the field $\F = R/\m$. By Corollary 11 of Hungerford \cite{hungerford} (\textbf{the use of this corollary is essentially the only new ingredient compared to the proof of Hillar and Rhea}), $R$ is the quotient of a PID which we call $Z$ (reminiscent of the role the integers $\Z$ play in the proof of Hillar and Rhea), and let $Q$ be the field of fractions of $Z$. (If $R$ itself is a domain, we may take $Z=R$.) Let $\hat{A} \in M_n(Z)$ be a lift of $A$, in the sense that when $\hat{A}_{ij}$ is taken modulo $\m^{e_i}$ for all $i$, we get back precisely $A_{ij}$. Let $S$ be the subset of $M_n(Z)$ whose image under the preceding map lies in $\overline{S}$; thus $\hat{A}$ is in $S$.

\begin{lemma}[{Cf. \cite[Lemma~3.2]{HillarRhea}}]
$S$ forms a ring under matrix multiplication. In particular, if we let $T$ be the diagonal matrix whose diagonal entries are $t^{e_1}, t^{e_2}, \dotsb, t^{e_n}$, then $A \in M_n(Z)$ is in $S$ if and only if there exists a matrix $A' \in M_n(Z)$ such that $A = TAT^{-1}$.
\end{lemma}

The proof of this lemma is almost immediate, and is identical, \emph{mutatis mutandis}, to that in Hillar \& Rhea.

By usual properties of the adjugate matrix (defined in terms of minors), there is a unique matrix $\hat{B} \in M_n(Z)$ such that $\hat{A}\hat{B} = \hat{B} \hat{A} = \det(\hat{A}) I_n$. We will show that $\hat{B}$ is in $S$, and thus its image in $\overline{S}$ is an inverse of $A$ in $\overline{S}$. 

By the preceding lemma, write $\hat{A} = TA' T^{-1}$, where $T$ is as in the lemma and $A' \in M_n(Z)$. Let $B'$ be the adjugate of $A'$. Let $B'' = TB'T^{-1}$. Then we have $\hat{A} B'' = T A' B' T^{-1} = \det(\hat{A}) I_n = T B' A' T^{-1} = B'' \hat{A}$. By the uniqueness of the adjugate, we have $B'' = \hat{B}$. Since $\hat{B}$ can thus be expressed as $T B' T^{-1}$ for some $B' \in M_n(Z)$, by the lemma we have that $\hat{B}$ is in $S$ as desired, and thus its image in $\overline{S}$ is an inverse of $A$, so $A$ is invertible.
\end{proof}

\begin{proof}[Proof of Lemma~\ref{lem:twisted}]
Given a tuple $\vec{B} = (B_1, \dotsc, B_m)$ of $\varepsilon$-symmetric bilinear forms over $M$, let $B_1, \dotsc, B_m$ also represent the corresponding matrices relative to some basis of $M$. If $M$ has a basis of size $n$, then let $M' \cong M \oplus M$, which has basis size $2n$. In a basis of $M'$ consisting of a basis of the first copy of $M$, followed by a basis of the second copy of $M$, let $\vec{B}' = (B_0', B_1', \dotsc, B_m')$ be the following tuple of $R$-module endomorphisms $M' \to M'$, viewing elements of $M'$ as column vectors of length $2n$:
\[
B_0' = \begin{pmatrix} \text{id}_M & 0 \\ 0 & 0 \end{pmatrix}
\qquad
B_i' = \begin{pmatrix} 0 & B_i \\ 0 & 0 \end{pmatrix} \quad (i=1,\dotsc,m).
\]
Similarly build a tuple $\vec{C}'$ from $\vec{C}$.

We claim that there exists $A,D \in \aut(M)$ such that $A \vec{B} = \vec{C}D$ if and only if there is $E \in \aut(M')$ such that $E \vec{B}' E^{-1} = \vec{C}'$, or equivalently $E \vec{B} = \vec{C} E$, or equivalently, the module structures on $M'$ given by $\vec{B}'$ and $\vec{C}'$ are isomorphic. To see this, write $E = \begin{pmatrix} A & G \\ H & D \end{pmatrix}$, where each of $A,G,H,D$ is an $R$-module homomorphism $M \to M$. From $E B_0' = C_0' E$, we get 
\[
\begin{pmatrix} A & 0 \\ H & 0 \end{pmatrix} = \begin{pmatrix} A & G \\ 0 & 0 \end{pmatrix},
\]
and therefore we have $G = H = 0$. Since $E$ is an invertible block-diagonal matrix, each diagonal block must be invertible, hence we have $A,D \in \aut(M)$. Furthermore, from $E B_i' = C_i' E$, we immediately get $A B_i = C_i D$. Thus $\vec{B}$ and $\vec{C}$ are twisted equivalent via the pair $(A^t, D)$. 

Conversely, if $A \vec{B} = \vec{C} D$ with $A,D$ invertible, then it is readily verified that taking $E = \begin{pmatrix} A & 0 \\ 0 & D \end{pmatrix}$ results in $A \vec{B}' = \vec{C}'$, and that $E$ is invertible. This completes the proof of the reduction.

The module isomorphism problem is then solved in polynomial time over the finite ring $R$ by the algorithm of Ciocanea-Teodorescu \cite{CT}.
\end{proof}